\section{Dérivées régularisées et valeurs spectrales impaires}
\label{ix-add:bendersky}

La continuation et l'équation fonctionnelle démontrées dans la
proposition~\ref{rz-num-03-2} permettent de relier deux opérations
sur la même fonctionnelle $Z$ : la dérivation aux entiers négatifs et l'
évaluation convergente aux entiers positifs impairs. Les coordonnées
régionales et l'opérateur de numération maintiennent leur généalogie antérieure.
La notation $\pi$ désigne la coordonnée HMT qui est déjà intervenue dans la
construction thêta--Mellin.

Soient $H_k=\sum_{j=1}^k1/j$, avec $H_0=0$, et $\mathsf B_j$ les nombres de Bernoulli construits par la récurrence de la section~\ref{lectura:manuscrito-04}, avec $\mathsf B_1=-1/2$. La typographie $\mathsf B_j$ distingue ces coefficients de la constante de Meissel--Mertens. On définit la hiérarchie de Bendersky--Glaisher--Kinkelin par
\[
 A_k=\exp\!\left(\frac{H_k\mathsf B_{k+1}}{k+1}-Z'(-k)\right),
 \qquad k\in\mathbb N_0.
\]
Les points $-k$ sont réguliers. Le terme de Bernoulli fixe la
normalisation correspondant à chaque niveau avant son évaluation.

\subsection{Covariance des dérivées régularisées}
On utilise ici la fonctionnelle $Z_\lambda=\lambda^{-s}Z$, dont le résidu
est $\lambda^{-1}$, et on définit
\[
 D_k(\lambda)=\exp[-Z_\lambda'(-k)].
\]
Cet objet conserve la normalisation de l'opérateur mis à l'échelle, distincte de la normalisation à résidu unitaire employée pour $\widehat Z_\lambda$ dans la section~\ref{ix-add:stieltjes}.

\HMTresultado{Proposition}{Transport des dérivées régularisées}{ix-add:regularized-scaling}
Pour $\lambda>0$ et $k\in\mathbb N_0$,
\[
 Z_\lambda'(-k)=\lambda^k
       \bigl(Z'(-k)-(\log\lambda)Z(-k)\bigr),
\]
\[
 D_k(\lambda)=D_k(1)^{\lambda^k}
               \lambda^{\lambda^k Z(-k)},\qquad
 D_k(1)=A_k\exp\!\left(-\frac{H_k\mathsf B_{k+1}}{k+1}\right).
\]
\textbf{Démonstration.} La règle du produit appliquée à
$\lambda^{-s}Z(s)$ donne la première identité. Son exponentielle négative
donne la seconde, et la définition de $A_k$ donne la troisième.
\hfill$\square$

Aux niveaux initiaux,
\[
\begin{aligned}
 A_0&=e^{-Z'(0)},\\
 \log A_1&=\frac1{12}-Z'(-1),\qquad \log A_2=-Z'(-2),\\
 \log A_3&=-\frac{11}{720}-Z'(-3),\qquad \log A_4=-Z'(-4).
\end{aligned}
\]
Par exemple, $D_1(1)=A_1e^{-1/12}$. Maintenir le normalisateur évite
d'identifier deux lecteurs différents parce qu'ils partagent la même dérivée.

\Needspace{12\baselineskip}
\subsection{Reconstruction des valeurs impaires}
\HMTresultado{Théorème}{Correspondance entre niveaux négatifs pairs et positifs impairs}{ix-add:odd-values}
Pour tout entier $n\geq1$,
\[
 Z'(-2n)=(-1)^n\frac{(2n)!}{2(2\pi)^{2n}}Z(2n+1),
\]
\[
 \frac{\log D_{2n}(\lambda)}{\lambda^{2n}}
 =\log A_{2n}
 =(-1)^{n+1}\frac{(2n)!}{2(2\pi)^{2n}}Z(2n+1).
\]
\textbf{Démonstration.} L'équation fonctionnelle de la forme complétée
implique l'identité méromorphe
\[
 \pi^{-s/2}\Gamma(s/2)Z(s)
 =\pi^{-(1-s)/2}\Gamma((1-s)/2)Z(1-s).
\]
En $s=-2n$, le second membre est régulier et le facteur gamma du
premier possède un pôle simple. Par conséquent, $Z(-2n)=0$. La récurrence
de gamma et son résidu en zéro donnent
\[
 \Gamma(s/2)=\frac{2(-1)^n}{n!(s+2n)}+O(1).
\]
De plus,
\[
 \Gamma\!\left(n+\frac12\right)
 =\frac{(2n)!\sqrt\pi}{4^nn!}.
\]
La comparaison des termes constants dans l'identité méromorphe produit
\[
 \pi^n\frac{2(-1)^n}{n!}Z'(-2n)
 =\pi^{-n}\frac{(2n)!}{4^nn!}Z(2n+1),
\]
qui est la première formule. La série
$t/(e^t-1)+t/2$ est paire ; d'où $\mathsf B_{2n+1}=0$ pour
$n\geq1$. La définition de $A_{2n}$ et la proposition précédente
achèvent la preuve. \hfill$\square$

En particulier, le premier niveau pair positif satisfait
\[
 Z(3)=4\pi^2\log A_2.
\]
La constante d'Apéry évaluée avec reste dans la section~\ref{lectura:manuscrito-04} est ainsi reliée à une dérivée régularisée de la même fonctionnelle. Les intervalles qui y sont démontrés sont transportés par l'exponentielle monotone une fois conservées les bornes de la coordonnée régionale $\pi$.

Le niveau zéro possède une évaluation différente. L'équation fonctionnelle,
conjointement avec la réflexion et la duplication de gamma, équivaut à
\[
 Z(s)=2^s\pi^{s-1}\sin(\pi s/2)\Gamma(1-s)Z(1-s).
\]
La formule d'Euler
\[
 \Gamma(1+z)=\lim_{m\to\infty}
 m^z\prod_{j=1}^{m}(1+z/j)^{-1}
\]
donne $\log\Gamma(1+z)=-\gamma_0z+O(z^2)$ : le coefficient linéaire est
la limite de $\log m-H_m$, et le reste quadratique est contrôlé
par la série convergente de $j^{-2}$. En utilisant
$Z(1-s)=-1/s+\gamma_0+O(s)$ et
$\Gamma(1-s)=1+\gamma_0s+O(s^2)$, les termes de degré un
qui contiennent $\gamma_0$ s'annulent et l'on obtient
\[
 Z(s)=-\frac12-\frac12\log(2\pi)s+O(s^2).
\]
En conséquence,
\[
 A_0=\sqrt{2\pi},\qquad D_0(\lambda)=\sqrt{\frac{2\pi}{\lambda}}.
\]
La condition $n\geq1$ du théorème conserve cette différence : le niveau
zéro ne s'obtient pas en substituant $Z(1)$, qui possède un pôle, dans la formule
des valeurs impaires. Les identités de gamma expriment la normalisation
analytique ultérieure, sans intervenir dans la sélection des germes.
